% Generated by scripts/assemble.py; corrections belong in source/corrections.json.
\examhead{123}{山东大学 2026 年数学分析试卷}{117}
% p118-b26; source PDF page 118
\blockanchor{p118-b26}
\question{1}求极限 $\lim_{n\to \infty}n\left(e - \left(1 + \frac{1}{n}\right)^n\right)$.\par
% p118-b27; source PDF page 118
\blockanchor{p118-b27}
\question{2}设 $f(x) = x^{\alpha}$, $x \in [1, +\infty)$, 求 $\alpha$ 的取值范围使得 $f(x)$ 在 $[1, +\infty)$ 上一致连续.\par
% p119-b00; source PDF page 119
\par\Needspace{6\baselineskip}
\blockanchor{p119-b00}
\question{3}已知定义在 $D = \{(x, y) \mid x^2 + y^2 \leq 1\}$ 上的二阶连续可微函数 $f(x, y)$ 满足 $\frac{\partial^2 f}{\partial x^2} + \frac{\partial^2 f}{\partial y^2} = x^2 + y^2$. 计算\par
\nopagebreak[4]
% p119-b01; source PDF page 119
\blockanchor{p119-b01}
\displayformula{\iint_ {D} \left(\frac {x}{\sqrt {x ^ {2} + y ^ {2}}} \frac {\partial f}{\partial x} + \frac {y}{\sqrt {x ^ {2} + y ^ {2}}} \frac {\partial f}{\partial y}\right) \mathrm{d} x \mathrm{d} y.}
% p119-b02; source PDF page 119
\blockanchor{p119-b02}
\question{4}已知 k > 0, 若  $\lim_{n\to\infty}\frac{b_{1}+b_{2}+\cdots+b_{n}}{n^{k}}=b$ , 求证:  $\lim_{n\to\infty}\frac{b_{n}}{n^{k}}=0$ .\par
% p119-b03; source PDF page 119
\blockanchor{p119-b03}
\question{5}已知 $f(x) \in C[0, 2\pi]$，证明：$\lim_{n \to \infty} \int_{0}^{2\pi} f(x) |\sin nx| \, \mathrm{d}x = \frac{2}{\pi} \int_{0}^{2\pi} f(x) \, \mathrm{d}x$.\par
% p119-b04; source PDF page 119
\blockanchor{p119-b04}
\question{6}设 $f(x) = \sum_{n=1}^{\infty} \frac{x^n}{n^2 \ln(1+n)}$，讨论 $f(x)$ 在 $[-1, 1]$ 上的连续性，以及在 $x = -1$ 和 $x = 1$ 时的可微性.\par
% p119-b05; source PDF page 119
\blockanchor{p119-b05}
\question{7}设 $f(x) = x^3 \left(e^{\frac{1}{x}} + e^{-\frac{1}{x}} - c\right)$ 满足 $\lim_{x \to \infty} [f(x) - (ax + b)] = 0$, 求 $a, b, c$, 并说明其几何意义.\par
% p119-b06; source PDF page 119
\par\Needspace{6\baselineskip}
\blockanchor{p119-b06}
\question{8}设 $f(t)$ 在 $[0, +\infty)$ 上连续，且\par
\nopagebreak[4]
% p119-b07; source PDF page 119
\blockanchor{p119-b07}
\displayformula{f (t) = \iint_ {x ^ {2} + y ^ {2} \leq t ^ {2}} \left(x ^ {2} + y ^ {2}\right) f \left(\sqrt {x ^ {2} + y ^ {2}}\right) \mathrm{d} x \mathrm{d} y + t ^ {4}, t > 0.}
% p119-b08; source PDF page 119
\blockanchor{p119-b08}
求 $f(t)$ .\par
% p119-b09; source PDF page 119
\blockanchor{p119-b09}
\question{9}已知球面 $\Gamma: x^{2} + y^{2} + z^{2} = a^{2} (a > 0)$, 球面 $\Sigma$ 的球心在 $\Gamma$ 上, 半径为 $R$, 当 $R$ 为何值时, 球面 $\Sigma$ 位于 $\Gamma$ 内的面积取最大值? 并求出最大面积.\par
